\documentclass[../../../include/open-logic-section]{subfiles}

\begin{document}

\olfileid[id]{sth}{choice}{banach}
\olsection{Paradoks Banach--Tarski}

Kita juga dapat mencoba membenarkan Pilihan, sebagaimana Boolos mencoba
membenarkan Penggantian, dengan mengacu pada pertimbangan
\emph{ekstrinsik} (lihat \olref[replacement][extrinsic]{sec}). Lagi pula,
menerima Pilihan mempunyai banyak konsekuensi yang diinginkan: kemampuan
membandingkan setiap kardinal; kemampuan mengurutkan baik setiap himpunan;
kemampuan memperlakukan kardinal sebagai suatu jenis ordinal tertentu; dan
sebagainya.

Namun, kadang-kadang diklaim bahwa Pilihan mempunyai konsekuensi yang
\emph{tidak diinginkan}. Sebagian besar klaim ini disebabkan oleh suatu
hasil dari \cite{BanachTarski1924}.

\begin{thm}[Paradoks Banach--Tarski (dalam $\ZFC$)]
Sebarang bola dapat diuraikan menjadi sejumlah hingga kepingan yang dapat
disusun kembali (melalui rotasi dan translasi) untuk membentuk dua salinan
bola tersebut.
\end{thm}
\noindent
Sekilas, hasil ini cukup mencengangkan. Jelas bahwa kedua bola mempunyai
\emph{dua kali} volume bola semula. Namun, gerak kaku---rotasi dan
translasi---tidak mengubah volume. Jadi, tampaknya Banach--Tarski
memungkinkan kita mengada-adakan materi baru.

Keadaannya bahkan lebih buruk.\footnote{Lihat
\citet[Teorema 3.12]{Wagon2016}.} Penalaran yang serupa menunjukkan bahwa
sebutir kacang polong dapat dipotong menjadi sejumlah hingga kepingan yang
kemudian dapat disusun kembali (melalui rotasi dan translasi) untuk membentuk
suatu benda dengan bentuk dan ukuran Big Ben.

Namun, tidak satu pun dari hasil ini berlaku hanya dalam $\ZF$.\footnote{
Meskipun Banach--Tarski dapat dibuktikan dengan prinsip-prinsip yang secara
ketat lebih lemah daripada Pilihan; lihat \citet[hlm.~303]{Wagon2016}.}
Jadi, kita menghadapi suatu keputusan: menolak Pilihan atau belajar hidup
dengan ``paradoks'' tersebut.

Kita akan mengemukakan bahwa kita sebaiknya belajar hidup dengan
``paradoks'' itu. Bahkan, kita tidak menganggapnya sebagai paradoks yang
berarti. Secara khusus, kita tidak melihat mengapa hasil itu lebih atau
kurang paradoksal daripada hasil-hasil berikut:\footnote{
\citet[hlm.~276--277]{Potter2004}, \citet[hlm.~16]{Weston2003}, dan
\citet[hlm.~31, 308--309]{Wagon2016} mengemukakan pokok serupa dengan
menggunakan contoh-contoh lain.}
\begin{enumerate}
	\item Terdapat sama banyak titik dalam interval $(0,1)$ dengan dalam
	$\Real$.
	\\\emph{Bukti}: perhatikan $\tan(\pi(r-\nicefrac{1}{2}))$.
	\item Terdapat sama banyak titik pada suatu garis dengan dalam suatu
	persegi.
	\\Lihat \olref[his][set][pathology]{sec} dan
	\olref[his][set][cantorplane]{sec}.
	\item Terdapat kurva-kurva pengisi ruang.
	\\Lihat \olref[his][set][pathology]{sec} dan
	\olref[his][set][hilbertcurve]{sec}.
\end{enumerate}
Tidak satu pun dari ketiga hasil ini memerlukan Pilihan. Bahkan, sekarang
kita cukup menganggapnya sebagai bagian matematika yang mengejutkan dan
indah. Barangkali kita sebaiknya bersikap sama terhadap Paradoks
Banach--Tarski.

Memang, suatu pengamatan teknis diperlukan di sini; tetapi kita hanya perlu
tetap berpikir jernih. Gerak kaku mempertahankan volume. Akibatnya, kelima
kepingan\footnote{Kita menyatakan Paradoks tersebut dalam istilah
``sejumlah hingga kepingan''. Sesungguhnya, \citet{Robinson1947}
membuktikan bahwa penguraian dapat dicapai dengan \emph{lima} kepingan
(tetapi tidak kurang). Untuk suatu bukti, lihat
\citet[hlm.~66--67]{Wagon2016}.} yang menjadi hasil penguraian bola tidak
mungkin semuanya \emph{terukur}. Secara kasar, tidak masuk akal untuk
menetapkan volume kepada kepingan-kepingan individual ini. Anda perlu
membayangkannya sebagai ``hamburan-hamburan tak hingga'' titik yang tidak
dapat digambarkan. Barangkali memang ``aneh'' untuk memikirkan himpunan
``yang terhambur secara tak hingga'' seperti itu. Namun, keberadaannya
tampak mengikuti perintah yang terkandung dalam \stagesacc{} bahwa kita
harus membentuk \emph{semua koleksi yang mungkin} dari himpunan-himpunan
yang telah tersedia sebelumnya.

Jika semua itu tidak meyakinkan, berikut suatu argumen (ekstrinsik) terakhir
untuk menerima Paradoks Banach--Tarski. Paradoks itu segera menghasilkan
lelucon matematika terbaik sepanjang masa:
\begin{enumerate}
	\item[] \emph{Pertanyaan}. Apa anagram dari ``Banach--Tarski''?
	\item[] \emph{Jawaban}. ``Banach--Tarski Banach--Tarski''.
\end{enumerate}

\end{document}
